\documentclass[12pt]{article}

\usepackage{fontawesome}
\usepackage{hyperref}
\usepackage{xurl}
\usepackage[tagged, highstructure]{accessibility}
\usepackage{bookmark}
\usepackage{enumitem}

\hypersetup{
    colorlinks=false,
    pdfborder={0 0 0},
    pdftitle={Containers Lab},
    pdfauthor={Adrianna Holden-Gouveia},
    pdfsubject={Linux Administration - Lab Assignment},
    pdflang={en-US},
    pdfkeywords={containers, Docker, Kubernetes, virtualization, server administration}
}

% Configure PDF bookmarks for navigation
\bookmarksetup{
    numbered,
    open,
}

% Configure list spacing for better accessibility
\setlist{nosep}


\title{Containers}
\author{
        Adrianna Holden-Gouveia \\
        Website: \href{https://aholdengouveia.name}{aholdengouveia.name}\\ 
        \date{\vspace{-5ex}}
        %Email: \href{mailto:admin@aholdengouveia.name}{admin@aholdengouveia.name} \\
%        \faLinkedin{: aholdengouveia} \\
        \faGithub {: aholdengouveia} \\
       % \faTwitter {: aholdengouveia} \\
        }

%S\date{\today}


\begin{document}    

\maketitle

%\begin{abstract}
%This is a template for Linux Administration Lab work
%\end{abstract}
%\tableofcontents


\section*{Accessibility Notice}
This document is also available in HTML format at:

Link: \href{https://aholdengouveia.name/LinuxAdmin/labexcercises/containers.html}{\textbf{\nolinkurl{https://aholdengouveia.name/LinuxAdmin/labexcercises/containers.html}}}

The HTML version provides enhanced accessibility features including keyboard navigation, screen reader support, responsive design, dark mode support, and high contrast options.


\section*{Ticket}
\noindent\textbf{Ticket:} ACME-8031\\
\textbf{Requested by:} IT Manager\\
\textbf{Systems:} One of your servers

\section*{Background}
Acme Corp is considering containers for future deployments, and your manager wants a proof of concept and a recommendation. You may do this on one of the two servers you set up in the Server Setup lab.

\section*{Learning Outcomes}
By the end of this lab, you will be able to:
\begin{itemize}
    \item Install and configure Docker and Kubernetes
    \item Deploy a basic web server in a container
    \item Evaluate the pros and cons of containers, and recommend whether a business should use them
\end{itemize}

\section*{Requirements}
\begin{enumerate}
    \item Install Docker and Kubernetes on your server. You can use MicroK8s as your Kubernetes installation method.
    \item Go through the minikube tutorial.
    \item Get a basic web server working using Docker or Kubernetes, your choice. You can use the Apache (httpd) Docker image, or if you would prefer, you can substitute another server such as Home Assistant or ZoneMinder. Take a screenshot showing your web server running and its page loading.
    \item Write an executive summary on what containers are, the pros and cons of containers, and why you think Acme Corp should or shouldn't consider using them going forward.
    \item Find at least one resource for containers that you found helpful and that I haven't already shared anywhere.
\end{enumerate}

\section*{Acceptance Criteria}
\begin{itemize}
    \item If you had to look anything up to figure out a command, cite your source with a link to the reference.
    \item The install process is documented with screenshots, explanations of what you're doing, and any troubleshooting that was needed.
    \item The web server is proven to work with a screenshot.
    \item The executive summary is less than 1 page.
    \item Your resource is new, not something I've already shared.
    \item Every terminal screenshot shows your prompt with your username and hostname, and either the output of \texttt{date} or a note beside the terminal with your name and the term. Hostnames must match your server inventory from the Server Setup lab.
\end{itemize}

\section*{Deliverables}
Audience: someone technically inclined but unfamiliar with virtualization. Include screenshots to walk them through the install process.
\begin{enumerate}
    \item A document on how to install Docker and Kubernetes on your server.
    \item A short document describing how the minikube tutorial went, what gave you trouble (if anything), and how you fixed it. Include screenshots showing the steps you did and minikube up and running.
    \item A screenshot of your working web server, with a short explanation of how you set it up.
    \item Your executive summary.
    \item A link to your container resource.
\end{enumerate}

\section*{Resources}
References, a video, a PowerPoint and some notes are available at my website
Link: \href{https://www.aholdengouveia.name/LinuxAdmin/containers.html}{\textbf{\nolinkurl{https://www.aholdengouveia.name/LinuxAdmin/containers.html}}}

MicroK8s installation instructions
Link: \href{https://microk8s.io/}{\textbf{\nolinkurl{https://microk8s.io/}}}

The minikube tutorial
Link: \href{https://kubernetes.io/docs/tutorials/hello-minikube/}{\textbf{\nolinkurl{https://kubernetes.io/docs/tutorials/hello-minikube/}}}

The Apache (httpd) Docker image
Link: \href{https://hub.docker.com/_/httpd}{\textbf{\nolinkurl{https://hub.docker.com/_/httpd}}}

My instructions and sample yaml files
Link: \href{https://www.aholdengouveia.name/SmartHome/Virtualizationsetup.html}{\textbf{\nolinkurl{https://www.aholdengouveia.name/SmartHome/Virtualizationsetup.html}}}

\end{document}